This chapter describes the implementation of SentAI as a research
platform positioned at the intersection of embedded AI acceleration,
real-time systems, and autonomous drone control. The central thesis of
the implementation is that an EdgeTPU-class accelerator is
simultaneously low-power and computationally powerful enough to justify
a much tighter integration with an MCU than is common in larger
companion-computer architectures. Rather than treating the accelerator
as a peripheral attached to a heavy host, the design treats the MCU and
EdgeTPU as a compact co-processing pair: the MCU handles sensing,
scheduling, memory orchestration, communication, and high-level control,
while the EdgeTPU is driven as the main execution substrate for neural
inference.

An equally important framing point is that the implementation starts
from an open-source platform with significant unrealized potential and
deliberately pushes it toward industrialization. In other words, the
work is not only about extending Coral Micro with additional
capabilities, but about transforming an accessible open-source base into
a system that exhibits more production-oriented behavior: higher
sustained throughput, better control of memory and latency, richer
sensor integration, clearer runtime abstractions, and a software
structure that can support repeated deployment rather than isolated
demonstrations.

The implementation presented in this work differs substantially from the
initial fork baseline by transforming a board-oriented demo environment
into a programmable autonomous-flight runtime. The main engineering
question is not only whether inference can run onboard, but how to
maximize accelerator utilization, sustain real-time responsiveness, and
preserve a lightweight and low-power hardware profile suitable for
micro-UAV operation. This leads to a design in which memory placement,
transfer scheduling, camera integration, scripting boundaries, and
telemetry abstractions become first-order implementation concerns. In
this sense, the chapter documents a move from open-source
proof-of-potential to an embedded system engineered with
production-performance objectives in mind.

The chapter is organized into the following subchapters:

\begin{enumerate}
\def\labelenumi{\arabic{enumi}.}
\tightlist
\item
  \url{implementation_01_system_architecture.md}, which explains the
  hardware-software co-design and the motivation for pairing an MCU with
  an EdgeTPU on a lightweight drone-oriented board.
\item
  \url{implementation_02_runtime_and_memory.md}, which analyzes the RTOS
  integration, linker-level memory layout, staging buffers, and the
  host-side logic required to keep the accelerator fed efficiently.
\item
  \url{implementation_03_vision_pipeline.md}, which details the
  dual-camera subsystem, the real-time image path, tracking, projection,
  and the deployment of YOLO-class detectors.
\item
  \url{implementation_04_flight_and_telemetry.md}, which presents the
  dual-target control strategy spanning Crazyflie and PX4/MAVLink,
  together with telemetry and backend integration.
\item
  \url{implementation_05_programming_model.md}, which studies the
  MicroPython layer and uses the platform help system as evidence of how
  the embedded capabilities are exposed as a coherent experimental
  interface.
\end{enumerate}

Taken together, these sections argue that the implementation is best
understood as an exercise in practical embedded-AI industrialization:
the work does not merely deploy neural inference on a small board, but
systematically resolves the hardware and software bottlenecks that arise
when such inference must coexist with flight control, camera switching,
runtime programmability, and constrained memory resources. The overall
contribution is therefore to show how an open-source platform can be
matured toward a more production-capable embedded autonomy stack without
abandoning its accessibility, compactness, or low-power design
philosophy.
